
\documentclass{article}
\usepackage{colortbl}
\usepackage{makecell}
\usepackage{multirow}
\usepackage{supertabular}

\begin{document}

\newcounter{utterance}

\centering \large Interaction Transcript for game `codenames', experiment `opponent\_difficult', episode 4 with deepseek{-}v4{-}pro.
\vspace{24pt}

{ \footnotesize  \setcounter{utterance}{1}
\setlength{\tabcolsep}{0pt}
\begin{supertabular}{c@{$\;$}|p{.15\linewidth}@{}p{.15\linewidth}p{.15\linewidth}p{.15\linewidth}p{.15\linewidth}p{.15\linewidth}}
    \# & $\;$A & \multicolumn{4}{c}{Game Master} & $\;\:$B\\
    \hline

    \theutterance \stepcounter{utterance}  
    & & & \multicolumn{2}{p{0.3\linewidth}}{
        \cellcolor[rgb]{0.9,0.9,0.9}{
            \makecell[{{p{\linewidth}}}]{
                \texttt{\tiny{[GM$|$GM]}}
                {'hidden': {'team': ['казино', 'лук', 'маса', 'пляшка', 'акції', 'англія', 'собака', 'ґудзик', 'день'], 'innocent': ['ангел', 'коло', 'амазонка', 'концерт', 'вчитель', 'стовп', 'америка'], 'opponent': ['вентилятор', 'час', 'лінійка', 'лимон', "зв'язок", 'муха', 'двір', "п'єса"], 'assassin': ['столиця']}, 'revealed': {'team': {'team': [], 'innocent': [], 'opponent': [], 'assassin': []}, 'opponent': {'team': [], 'innocent': [], 'opponent': [], 'assassin': []}}}
            }
        }
    }
    & & \\ \\

    \theutterance \stepcounter{utterance}  
    & & \multicolumn{4}{p{0.6\linewidth}}{
        \cellcolor[rgb]{0.9,0.9,0.9}{
            \makecell[{{p{\linewidth}}}]{
                \texttt{\tiny{[P1$\langle$GM]}}
                \texttt{Надайте одну однословну підказку, яка пов'язана з деякими словами вашої команди зі списку нижче. Ви можете обрати, зі скількома і з якими саме словами пов'язана ваша підказка.} \\
\texttt{Мета полягає в тому, щоб ваш партнер відгадав усі слова вашої команди на основі підказок, які ви йому даєте, раніше, ніж інша команда відкриє всі свої. Якщо ваш партнер відгадає слово «суперника» або «відволікача», його хід відгадування закінчується, і команда суперників отримує право зробити хід. Якщо ваш партнер відгадає слово «вбивці», гра закінчується, і ваша команда негайно програє.} \\
\texttt{Чим із більшою кількістю слів пов'язана ваша підказка, тим швидше ви можете виграти гру, але тим складнішими можуть бути асоціації для відгадування вашим партнером, тож обирайте слово{-}підказку мудро. Слово{-}підказка має бути семантично пов'язаним із цільовими словами; воно не може бути одним зі слів у списках або містити їхні частини.} \\
\\ 
\texttt{Завжди надавайте вашу однословну підказку та ваш розділений комами список пов'язаних цільових слів у такому форматі й робіть відповіді якомога коротшими; ніколи не додавайте жодного іншого тексту, окрім того, що потрібно в цій формі:} \\
\\ 
\texttt{ПІДКАЗКА: <СЛОВО>} \\
\texttt{ЦІЛІ: <ЦІЛІ>} \\
\\ 
\texttt{Слова вашої команди: казино, лук, маса, пляшка, акції, англія, собака, ґудзик, день.} \\
\texttt{Слова команди суперника: вентилятор, час, лінійка, лимон, зв'язок, муха, двір, п'єса.} \\
\texttt{Слова{-}відволікачі: ангел, коло, амазонка, концерт, вчитель, стовп, америка.} \\
\texttt{Слова{-}вбивці: столиця.} \\
            }
        }
    }
    & \\ \\

    \theutterance \stepcounter{utterance}  
    & \multicolumn{4}{p{0.6\linewidth}}{
        \cellcolor[rgb]{0.9,0.9,0.9}{
            \makecell[{{p{\linewidth}}}]{
                \texttt{\tiny{[P1$\rangle$GM]}}
                \texttt{ПІДКАЗКА: Робін} \\
\texttt{ЦІЛІ: лук, англія, собака, день} \\
            }
        }
    }
    & & \\ \\

    \theutterance \stepcounter{utterance}  
    & & & \multicolumn{2}{p{0.3\linewidth}}{
        \cellcolor[rgb]{0.9,0.9,0.9}{
            \makecell[{{p{\linewidth}}}]{
                \texttt{\tiny{[GM$|$GM]}}
                \texttt{робін} \\
            }
        }
    }
    & & \\ \\

    \theutterance \stepcounter{utterance}  
    & & & \multicolumn{2}{p{0.3\linewidth}}{
        \cellcolor[rgb]{0.9,0.9,0.9}{
            \makecell[{{p{\linewidth}}}]{
                \texttt{\tiny{[GM$|$GM]}}
                ['лук', 'англія', 'собака', 'день']
            }
        }
    }
    & & \\ \\

    \theutterance \stepcounter{utterance}  
    & & & \multicolumn{2}{p{0.3\linewidth}}{
        \cellcolor[rgb]{0.9,0.9,0.9}{
            \makecell[{{p{\linewidth}}}]{
                \texttt{\tiny{[GM$|$GM]}}
                {'word': 'лук', 'assignment': 'team'}
            }
        }
    }
    & & \\ \\

    \theutterance \stepcounter{utterance}  
    & & & \multicolumn{2}{p{0.3\linewidth}}{
        \cellcolor[rgb]{0.9,0.9,0.9}{
            \makecell[{{p{\linewidth}}}]{
                \texttt{\tiny{[GM$|$GM]}}
                {'word': 'англія', 'assignment': 'team'}
            }
        }
    }
    & & \\ \\

    \theutterance \stepcounter{utterance}  
    & & & \multicolumn{2}{p{0.3\linewidth}}{
        \cellcolor[rgb]{0.9,0.9,0.9}{
            \makecell[{{p{\linewidth}}}]{
                \texttt{\tiny{[GM$|$GM]}}
                {'word': 'собака', 'assignment': 'team'}
            }
        }
    }
    & & \\ \\

    \theutterance \stepcounter{utterance}  
    & & & \multicolumn{2}{p{0.3\linewidth}}{
        \cellcolor[rgb]{0.9,0.9,0.9}{
            \makecell[{{p{\linewidth}}}]{
                \texttt{\tiny{[GM$|$GM]}}
                {'word': 'день', 'assignment': 'team'}
            }
        }
    }
    & & \\ \\

    \theutterance \stepcounter{utterance}  
    & & \multicolumn{4}{p{0.6\linewidth}}{
        \cellcolor[rgb]{0.9,0.9,0.9}{
            \makecell[{{p{\linewidth}}}]{
                \texttt{\tiny{[GM$\rangle$P2]}}
                \texttt{Надайте розділений комами список із не більш ніж 4 слів із наведеного нижче списку, які найкраще пов'язані або найбільш тісно асоціюються зі словом «робін». Завжди починайте свій список здогадок із ВІДГАДКА: і не додавайте жодного іншого тексту у своїй відповіді.} \\
\\ 
\texttt{лінійка, лук, двір, ангел, америка, собака, день, вчитель, концерт, стовп, маса, англія, столиця, вентилятор, казино, муха, амазонка, зв'язок, пляшка, акції, час, ґудзик, коло, лимон, п'єса} \\
            }
        }
    }
    & \\ \\

    \theutterance \stepcounter{utterance}  
    & & & \multicolumn{4}{p{0.6\linewidth}}{
        \cellcolor[rgb]{0.9,0.9,0.9}{
            \makecell[{{p{\linewidth}}}]{
                \texttt{\tiny{[GM$\langle$P2]}}
                \texttt{ВІДГАДКА: лук, муха, англія, день} \\
            }
        }
    }
     \\ \\

    \theutterance \stepcounter{utterance}  
    & & & \multicolumn{2}{p{0.3\linewidth}}{
        \cellcolor[rgb]{0.9,0.9,0.9}{
            \makecell[{{p{\linewidth}}}]{
                \texttt{\tiny{[GM$|$GM]}}
                ['лук', 'муха', 'англія', 'день']
            }
        }
    }
    & & \\ \\

    \theutterance \stepcounter{utterance}  
    & & & \multicolumn{2}{p{0.3\linewidth}}{
        \cellcolor[rgb]{0.9,0.9,0.9}{
            \makecell[{{p{\linewidth}}}]{
                \texttt{\tiny{[GM$|$GM]}}
                {'word': 'лук', 'assignment': 'team'}
            }
        }
    }
    & & \\ \\

    \theutterance \stepcounter{utterance}  
    & & & \multicolumn{2}{p{0.3\linewidth}}{
        \cellcolor[rgb]{0.9,0.9,0.9}{
            \makecell[{{p{\linewidth}}}]{
                \texttt{\tiny{[GM$|$GM]}}
                {'word': 'лук', 'assignment': 'team'}
            }
        }
    }
    & & \\ \\

    \theutterance \stepcounter{utterance}  
    & & & \multicolumn{2}{p{0.3\linewidth}}{
        \cellcolor[rgb]{0.9,0.9,0.9}{
            \makecell[{{p{\linewidth}}}]{
                \texttt{\tiny{[GM$|$GM]}}
                {'word': 'муха', 'assignment': 'opponent'}
            }
        }
    }
    & & \\ \\

    \theutterance \stepcounter{utterance}  
    & & & \multicolumn{2}{p{0.3\linewidth}}{
        \cellcolor[rgb]{0.9,0.9,0.9}{
            \makecell[{{p{\linewidth}}}]{
                \texttt{\tiny{[GM$|$GM]}}
                \texttt{муха} \\
            }
        }
    }
    & & \\ \\

    \theutterance \stepcounter{utterance}  
    & & & \multicolumn{2}{p{0.3\linewidth}}{
        \cellcolor[rgb]{0.9,0.9,0.9}{
            \makecell[{{p{\linewidth}}}]{
                \texttt{\tiny{[GM$|$GM]}}
                {'word': 'час', 'assignment': 'opponent'}
            }
        }
    }
    & & \\ \\

    \theutterance \stepcounter{utterance}  
    & & & \multicolumn{2}{p{0.3\linewidth}}{
        \cellcolor[rgb]{0.9,0.9,0.9}{
            \makecell[{{p{\linewidth}}}]{
                \texttt{\tiny{[GM$|$GM]}}
                {'word': "зв'язок", 'assignment': 'opponent'}
            }
        }
    }
    & & \\ \\

    \theutterance \stepcounter{utterance}  
    & & & \multicolumn{2}{p{0.3\linewidth}}{
        \cellcolor[rgb]{0.9,0.9,0.9}{
            \makecell[{{p{\linewidth}}}]{
                \texttt{\tiny{[GM$|$GM]}}
                {'hidden': {'team': ['казино', 'маса', 'пляшка', 'акції', 'англія', 'собака', 'ґудзик', 'день'], 'innocent': ['ангел', 'коло', 'амазонка', 'концерт', 'вчитель', 'стовп', 'америка'], 'opponent': ['вентилятор', 'лінійка', 'лимон', 'двір', "п'єса"], 'assassin': ['столиця']}, 'revealed': {'team': {'team': ['лук'], 'innocent': [], 'opponent': ['муха'], 'assassin': []}, 'opponent': {'team': [], 'innocent': [], 'opponent': ['час', "зв'язок"], 'assassin': []}}}
            }
        }
    }
    & & \\ \\

    \theutterance \stepcounter{utterance}  
    & & \multicolumn{4}{p{0.6\linewidth}}{
        \cellcolor[rgb]{0.9,0.9,0.9}{
            \makecell[{{p{\linewidth}}}]{
                \texttt{\tiny{[P1$\langle$GM]}}
                \texttt{Слово лук було відгадано правильно. Слово муха було відгадано, але це суперницьке слово. Хід вашого товариша по команді закінчився.} \\
\texttt{Тепер надайте ще одну підказку, пов'язану з деякими з ваших командних слів, що залишилися, та список пов'язаних цільових слів. Пам'ятайте: почніть підказку з 'ПІДКАЗКА: ', перейдіть на новий рядок і почніть ваш список цільових слів, розділених комами, з 'ЦІЛІ: '. Зверніть увагу: деякі слова було вилучено зі списків порівняно з попередніми запитами.} \\
\\ 
\texttt{Ваші командні слова, що залишилися: казино, маса, пляшка, акції, англія, собака, ґудзик, день.} \\
\texttt{Слова, що залишилися для вашого суперника: вентилятор, час, лінійка, лимон, зв'язок, двір, п'єса.} \\
\texttt{Слова{-}відволікачі, що залишилися: ангел, коло, амазонка, концерт, вчитель, стовп, америка.} \\
\texttt{Слова{-}вбивці, що залишилися: столиця.} \\
            }
        }
    }
    & \\ \\

    \theutterance \stepcounter{utterance}  
    & \multicolumn{4}{p{0.6\linewidth}}{
        \cellcolor[rgb]{0.9,0.9,0.9}{
            \makecell[{{p{\linewidth}}}]{
                \texttt{\tiny{[P1$\rangle$GM]}}
                \texttt{ПІДКАЗКА: туманний} \\
\texttt{ЦІЛІ: англія, день} \\
            }
        }
    }
    & & \\ \\

    \theutterance \stepcounter{utterance}  
    & & & \multicolumn{2}{p{0.3\linewidth}}{
        \cellcolor[rgb]{0.9,0.9,0.9}{
            \makecell[{{p{\linewidth}}}]{
                \texttt{\tiny{[GM$|$GM]}}
                \texttt{туманний} \\
            }
        }
    }
    & & \\ \\

    \theutterance \stepcounter{utterance}  
    & & & \multicolumn{2}{p{0.3\linewidth}}{
        \cellcolor[rgb]{0.9,0.9,0.9}{
            \makecell[{{p{\linewidth}}}]{
                \texttt{\tiny{[GM$|$GM]}}
                ['англія', 'день']
            }
        }
    }
    & & \\ \\

    \theutterance \stepcounter{utterance}  
    & & & \multicolumn{2}{p{0.3\linewidth}}{
        \cellcolor[rgb]{0.9,0.9,0.9}{
            \makecell[{{p{\linewidth}}}]{
                \texttt{\tiny{[GM$|$GM]}}
                {'word': 'англія', 'assignment': 'team'}
            }
        }
    }
    & & \\ \\

    \theutterance \stepcounter{utterance}  
    & & & \multicolumn{2}{p{0.3\linewidth}}{
        \cellcolor[rgb]{0.9,0.9,0.9}{
            \makecell[{{p{\linewidth}}}]{
                \texttt{\tiny{[GM$|$GM]}}
                {'word': 'день', 'assignment': 'team'}
            }
        }
    }
    & & \\ \\

    \theutterance \stepcounter{utterance}  
    & & \multicolumn{4}{p{0.6\linewidth}}{
        \cellcolor[rgb]{0.9,0.9,0.9}{
            \makecell[{{p{\linewidth}}}]{
                \texttt{\tiny{[GM$\rangle$P2]}}
                \texttt{Слово лук було відгадано правильно. Слово муха було відгадано, але це суперницьке слово. Ваш хід закінчився.} \\
\texttt{Тепер надайте ще один список слів, розділених комами, щонайменше з 1 і щонайбільше з 2 слів із наведеного нижче списку слів, які найкраще пов'язані або найбільш тісно асоціюються зі словом 'туманний'. Пам'ятайте, що потрібно почати вашу відповідь з ВІДГАДКА:. Зауваження: деякі слова було вилучено зі списку порівняно з попередніми запитами.} \\
\\ 
\texttt{лінійка, двір, ангел, америка, собака, день, вчитель, концерт, стовп, маса, англія, столиця, вентилятор, казино, амазонка, пляшка, акції, ґудзик, коло, лимон, п'єса} \\
            }
        }
    }
    & \\ \\

    \theutterance \stepcounter{utterance}  
    & & & \multicolumn{4}{p{0.6\linewidth}}{
        \cellcolor[rgb]{0.9,0.9,0.9}{
            \makecell[{{p{\linewidth}}}]{
                \texttt{\tiny{[GM$\langle$P2]}}
                \texttt{ВІДГАДКА: день, англія} \\
            }
        }
    }
     \\ \\

    \theutterance \stepcounter{utterance}  
    & & & \multicolumn{2}{p{0.3\linewidth}}{
        \cellcolor[rgb]{0.9,0.9,0.9}{
            \makecell[{{p{\linewidth}}}]{
                \texttt{\tiny{[GM$|$GM]}}
                ['день', 'англія']
            }
        }
    }
    & & \\ \\

    \theutterance \stepcounter{utterance}  
    & & & \multicolumn{2}{p{0.3\linewidth}}{
        \cellcolor[rgb]{0.9,0.9,0.9}{
            \makecell[{{p{\linewidth}}}]{
                \texttt{\tiny{[GM$|$GM]}}
                {'word': 'день', 'assignment': 'team'}
            }
        }
    }
    & & \\ \\

    \theutterance \stepcounter{utterance}  
    & & & \multicolumn{2}{p{0.3\linewidth}}{
        \cellcolor[rgb]{0.9,0.9,0.9}{
            \makecell[{{p{\linewidth}}}]{
                \texttt{\tiny{[GM$|$GM]}}
                {'word': 'день', 'assignment': 'team'}
            }
        }
    }
    & & \\ \\

    \theutterance \stepcounter{utterance}  
    & & & \multicolumn{2}{p{0.3\linewidth}}{
        \cellcolor[rgb]{0.9,0.9,0.9}{
            \makecell[{{p{\linewidth}}}]{
                \texttt{\tiny{[GM$|$GM]}}
                {'word': 'англія', 'assignment': 'team'}
            }
        }
    }
    & & \\ \\

    \theutterance \stepcounter{utterance}  
    & & & \multicolumn{2}{p{0.3\linewidth}}{
        \cellcolor[rgb]{0.9,0.9,0.9}{
            \makecell[{{p{\linewidth}}}]{
                \texttt{\tiny{[GM$|$GM]}}
                {'word': 'англія', 'assignment': 'team'}
            }
        }
    }
    & & \\ \\

    \theutterance \stepcounter{utterance}  
    & & & \multicolumn{2}{p{0.3\linewidth}}{
        \cellcolor[rgb]{0.9,0.9,0.9}{
            \makecell[{{p{\linewidth}}}]{
                \texttt{\tiny{[GM$|$GM]}}
                {'word': 'двір', 'assignment': 'opponent'}
            }
        }
    }
    & & \\ \\

    \theutterance \stepcounter{utterance}  
    & & & \multicolumn{2}{p{0.3\linewidth}}{
        \cellcolor[rgb]{0.9,0.9,0.9}{
            \makecell[{{p{\linewidth}}}]{
                \texttt{\tiny{[GM$|$GM]}}
                {'word': 'вентилятор', 'assignment': 'opponent'}
            }
        }
    }
    & & \\ \\

    \theutterance \stepcounter{utterance}  
    & & & \multicolumn{2}{p{0.3\linewidth}}{
        \cellcolor[rgb]{0.9,0.9,0.9}{
            \makecell[{{p{\linewidth}}}]{
                \texttt{\tiny{[GM$|$GM]}}
                {'hidden': {'team': ['казино', 'маса', 'пляшка', 'акції', 'собака', 'ґудзик'], 'innocent': ['ангел', 'коло', 'амазонка', 'концерт', 'вчитель', 'стовп', 'америка'], 'opponent': ['лінійка', 'лимон', "п'єса"], 'assassin': ['столиця']}, 'revealed': {'team': {'team': ['лук', 'день', 'англія'], 'innocent': [], 'opponent': ['муха'], 'assassin': []}, 'opponent': {'team': [], 'innocent': [], 'opponent': ['час', "зв'язок", 'двір', 'вентилятор'], 'assassin': []}}}
            }
        }
    }
    & & \\ \\

    \theutterance \stepcounter{utterance}  
    & & \multicolumn{4}{p{0.6\linewidth}}{
        \cellcolor[rgb]{0.9,0.9,0.9}{
            \makecell[{{p{\linewidth}}}]{
                \texttt{\tiny{[P1$\langle$GM]}}
                \texttt{Слова день, англія було відгадано правильно. Хід вашого товариша по команді закінчився.} \\
\texttt{Тепер надайте ще одну підказку, пов'язану з деякими з ваших командних слів, що залишилися, та список пов'язаних цільових слів. Пам'ятайте: почніть підказку з 'ПІДКАЗКА: ', перейдіть на новий рядок і почніть ваш список цільових слів, розділених комами, з 'ЦІЛІ: '. Зверніть увагу: деякі слова було вилучено зі списків порівняно з попередніми запитами.} \\
\\ 
\texttt{Ваші командні слова, що залишилися: казино, маса, пляшка, акції, собака, ґудзик.} \\
\texttt{Слова, що залишилися для вашого суперника: вентилятор, лінійка, лимон, двір, п'єса.} \\
\texttt{Слова{-}відволікачі, що залишилися: ангел, коло, амазонка, концерт, вчитель, стовп, америка.} \\
\texttt{Слова{-}вбивці, що залишилися: столиця.} \\
            }
        }
    }
    & \\ \\

    \theutterance \stepcounter{utterance}  
    & \multicolumn{4}{p{0.6\linewidth}}{
        \cellcolor[rgb]{0.9,0.9,0.9}{
            \makecell[{{p{\linewidth}}}]{
                \texttt{\tiny{[P1$\rangle$GM]}}
                \texttt{ПІДКАЗКА: фішки} \\
\texttt{ЦІЛІ: казино, маса, акції} \\
            }
        }
    }
    & & \\ \\

    \theutterance \stepcounter{utterance}  
    & & & \multicolumn{2}{p{0.3\linewidth}}{
        \cellcolor[rgb]{0.9,0.9,0.9}{
            \makecell[{{p{\linewidth}}}]{
                \texttt{\tiny{[GM$|$GM]}}
                \texttt{фішки} \\
            }
        }
    }
    & & \\ \\

    \theutterance \stepcounter{utterance}  
    & & & \multicolumn{2}{p{0.3\linewidth}}{
        \cellcolor[rgb]{0.9,0.9,0.9}{
            \makecell[{{p{\linewidth}}}]{
                \texttt{\tiny{[GM$|$GM]}}
                ['казино', 'маса', 'акції']
            }
        }
    }
    & & \\ \\

    \theutterance \stepcounter{utterance}  
    & & & \multicolumn{2}{p{0.3\linewidth}}{
        \cellcolor[rgb]{0.9,0.9,0.9}{
            \makecell[{{p{\linewidth}}}]{
                \texttt{\tiny{[GM$|$GM]}}
                {'word': 'казино', 'assignment': 'team'}
            }
        }
    }
    & & \\ \\

    \theutterance \stepcounter{utterance}  
    & & & \multicolumn{2}{p{0.3\linewidth}}{
        \cellcolor[rgb]{0.9,0.9,0.9}{
            \makecell[{{p{\linewidth}}}]{
                \texttt{\tiny{[GM$|$GM]}}
                {'word': 'маса', 'assignment': 'team'}
            }
        }
    }
    & & \\ \\

    \theutterance \stepcounter{utterance}  
    & & & \multicolumn{2}{p{0.3\linewidth}}{
        \cellcolor[rgb]{0.9,0.9,0.9}{
            \makecell[{{p{\linewidth}}}]{
                \texttt{\tiny{[GM$|$GM]}}
                {'word': 'акції', 'assignment': 'team'}
            }
        }
    }
    & & \\ \\

    \theutterance \stepcounter{utterance}  
    & & \multicolumn{4}{p{0.6\linewidth}}{
        \cellcolor[rgb]{0.9,0.9,0.9}{
            \makecell[{{p{\linewidth}}}]{
                \texttt{\tiny{[GM$\rangle$P2]}}
                \texttt{Слова день, англія було відгадано правильно. Ваш хід закінчився.} \\
\texttt{Тепер надайте ще один список слів, розділених комами, щонайменше з 1 і щонайбільше з 3 слів із наведеного нижче списку слів, які найкраще пов'язані або найбільш тісно асоціюються зі словом 'фішки'. Пам'ятайте, що потрібно почати вашу відповідь з ВІДГАДКА:. Зауваження: деякі слова було вилучено зі списку порівняно з попередніми запитами.} \\
\\ 
\texttt{лінійка, ангел, америка, собака, вчитель, концерт, стовп, маса, столиця, казино, амазонка, пляшка, акції, ґудзик, коло, лимон, п'єса} \\
            }
        }
    }
    & \\ \\

    \theutterance \stepcounter{utterance}  
    & & & \multicolumn{4}{p{0.6\linewidth}}{
        \cellcolor[rgb]{0.9,0.9,0.9}{
            \makecell[{{p{\linewidth}}}]{
                \texttt{\tiny{[GM$\langle$P2]}}
                \texttt{ВІДГАДКА: казино, акції} \\
            }
        }
    }
     \\ \\

    \theutterance \stepcounter{utterance}  
    & & & \multicolumn{2}{p{0.3\linewidth}}{
        \cellcolor[rgb]{0.9,0.9,0.9}{
            \makecell[{{p{\linewidth}}}]{
                \texttt{\tiny{[GM$|$GM]}}
                ['казино', 'акції']
            }
        }
    }
    & & \\ \\

    \theutterance \stepcounter{utterance}  
    & & & \multicolumn{2}{p{0.3\linewidth}}{
        \cellcolor[rgb]{0.9,0.9,0.9}{
            \makecell[{{p{\linewidth}}}]{
                \texttt{\tiny{[GM$|$GM]}}
                {'word': 'казино', 'assignment': 'team'}
            }
        }
    }
    & & \\ \\

    \theutterance \stepcounter{utterance}  
    & & & \multicolumn{2}{p{0.3\linewidth}}{
        \cellcolor[rgb]{0.9,0.9,0.9}{
            \makecell[{{p{\linewidth}}}]{
                \texttt{\tiny{[GM$|$GM]}}
                {'word': 'казино', 'assignment': 'team'}
            }
        }
    }
    & & \\ \\

    \theutterance \stepcounter{utterance}  
    & & & \multicolumn{2}{p{0.3\linewidth}}{
        \cellcolor[rgb]{0.9,0.9,0.9}{
            \makecell[{{p{\linewidth}}}]{
                \texttt{\tiny{[GM$|$GM]}}
                {'word': 'акції', 'assignment': 'team'}
            }
        }
    }
    & & \\ \\

    \theutterance \stepcounter{utterance}  
    & & & \multicolumn{2}{p{0.3\linewidth}}{
        \cellcolor[rgb]{0.9,0.9,0.9}{
            \makecell[{{p{\linewidth}}}]{
                \texttt{\tiny{[GM$|$GM]}}
                {'word': 'акції', 'assignment': 'team'}
            }
        }
    }
    & & \\ \\

    \theutterance \stepcounter{utterance}  
    & & & \multicolumn{2}{p{0.3\linewidth}}{
        \cellcolor[rgb]{0.9,0.9,0.9}{
            \makecell[{{p{\linewidth}}}]{
                \texttt{\tiny{[GM$|$GM]}}
                {'word': 'лінійка', 'assignment': 'opponent'}
            }
        }
    }
    & & \\ \\

    \theutterance \stepcounter{utterance}  
    & & & \multicolumn{2}{p{0.3\linewidth}}{
        \cellcolor[rgb]{0.9,0.9,0.9}{
            \makecell[{{p{\linewidth}}}]{
                \texttt{\tiny{[GM$|$GM]}}
                {'word': "п'єса", 'assignment': 'opponent'}
            }
        }
    }
    & & \\ \\

    \theutterance \stepcounter{utterance}  
    & & & \multicolumn{2}{p{0.3\linewidth}}{
        \cellcolor[rgb]{0.9,0.9,0.9}{
            \makecell[{{p{\linewidth}}}]{
                \texttt{\tiny{[GM$|$GM]}}
                {'hidden': {'team': ['маса', 'пляшка', 'собака', 'ґудзик'], 'innocent': ['ангел', 'коло', 'амазонка', 'концерт', 'вчитель', 'стовп', 'америка'], 'opponent': ['лимон'], 'assassin': ['столиця']}, 'revealed': {'team': {'team': ['лук', 'день', 'англія', 'казино', 'акції'], 'innocent': [], 'opponent': ['муха'], 'assassin': []}, 'opponent': {'team': [], 'innocent': [], 'opponent': ['час', "зв'язок", 'двір', 'вентилятор', 'лінійка', "п'єса"], 'assassin': []}}}
            }
        }
    }
    & & \\ \\

    \theutterance \stepcounter{utterance}  
    & & \multicolumn{4}{p{0.6\linewidth}}{
        \cellcolor[rgb]{0.9,0.9,0.9}{
            \makecell[{{p{\linewidth}}}]{
                \texttt{\tiny{[P1$\langle$GM]}}
                \texttt{Слова казино, акції було відгадано правильно. Хід вашого товариша по команді закінчився.} \\
\texttt{Тепер надайте ще одну підказку, пов'язану з деякими з ваших командних слів, що залишилися, та список пов'язаних цільових слів. Пам'ятайте: почніть підказку з 'ПІДКАЗКА: ', перейдіть на новий рядок і почніть ваш список цільових слів, розділених комами, з 'ЦІЛІ: '. Зверніть увагу: деякі слова було вилучено зі списків порівняно з попередніми запитами.} \\
\\ 
\texttt{Ваші командні слова, що залишилися: маса, пляшка, собака, ґудзик.} \\
\texttt{Слова, що залишилися для вашого суперника: лінійка, лимон, п'єса.} \\
\texttt{Слова{-}відволікачі, що залишилися: ангел, коло, амазонка, концерт, вчитель, стовп, америка.} \\
\texttt{Слова{-}вбивці, що залишилися: столиця.} \\
            }
        }
    }
    & \\ \\

    \theutterance \stepcounter{utterance}  
    & \multicolumn{4}{p{0.6\linewidth}}{
        \cellcolor[rgb]{0.9,0.9,0.9}{
            \makecell[{{p{\linewidth}}}]{
                \texttt{\tiny{[P1$\rangle$GM]}}
                \texttt{ПІДКАЗКА: застібка} \\
\texttt{ЦІЛІ: ґудзик, собака} \\
            }
        }
    }
    & & \\ \\

    \theutterance \stepcounter{utterance}  
    & & & \multicolumn{2}{p{0.3\linewidth}}{
        \cellcolor[rgb]{0.9,0.9,0.9}{
            \makecell[{{p{\linewidth}}}]{
                \texttt{\tiny{[GM$|$GM]}}
                \texttt{застібка} \\
            }
        }
    }
    & & \\ \\

    \theutterance \stepcounter{utterance}  
    & & & \multicolumn{2}{p{0.3\linewidth}}{
        \cellcolor[rgb]{0.9,0.9,0.9}{
            \makecell[{{p{\linewidth}}}]{
                \texttt{\tiny{[GM$|$GM]}}
                ['ґудзик', 'собака']
            }
        }
    }
    & & \\ \\

    \theutterance \stepcounter{utterance}  
    & & & \multicolumn{2}{p{0.3\linewidth}}{
        \cellcolor[rgb]{0.9,0.9,0.9}{
            \makecell[{{p{\linewidth}}}]{
                \texttt{\tiny{[GM$|$GM]}}
                {'word': 'ґудзик', 'assignment': 'team'}
            }
        }
    }
    & & \\ \\

    \theutterance \stepcounter{utterance}  
    & & & \multicolumn{2}{p{0.3\linewidth}}{
        \cellcolor[rgb]{0.9,0.9,0.9}{
            \makecell[{{p{\linewidth}}}]{
                \texttt{\tiny{[GM$|$GM]}}
                {'word': 'собака', 'assignment': 'team'}
            }
        }
    }
    & & \\ \\

    \theutterance \stepcounter{utterance}  
    & & \multicolumn{4}{p{0.6\linewidth}}{
        \cellcolor[rgb]{0.9,0.9,0.9}{
            \makecell[{{p{\linewidth}}}]{
                \texttt{\tiny{[GM$\rangle$P2]}}
                \texttt{Слова казино, акції було відгадано правильно. Ваш хід закінчився.} \\
\texttt{Тепер надайте ще один список слів, розділених комами, щонайменше з 1 і щонайбільше з 2 слів із наведеного нижче списку слів, які найкраще пов'язані або найбільш тісно асоціюються зі словом 'застібка'. Пам'ятайте, що потрібно почати вашу відповідь з ВІДГАДКА:. Зауваження: деякі слова було вилучено зі списку порівняно з попередніми запитами.} \\
\\ 
\texttt{ангел, америка, собака, вчитель, концерт, стовп, маса, столиця, амазонка, пляшка, ґудзик, коло, лимон} \\
            }
        }
    }
    & \\ \\

    \theutterance \stepcounter{utterance}  
    & & & \multicolumn{4}{p{0.6\linewidth}}{
        \cellcolor[rgb]{0.9,0.9,0.9}{
            \makecell[{{p{\linewidth}}}]{
                \texttt{\tiny{[GM$\langle$P2]}}
                \texttt{ВІДГАДКА: ґудзик} \\
            }
        }
    }
     \\ \\

    \theutterance \stepcounter{utterance}  
    & & & \multicolumn{2}{p{0.3\linewidth}}{
        \cellcolor[rgb]{0.9,0.9,0.9}{
            \makecell[{{p{\linewidth}}}]{
                \texttt{\tiny{[GM$|$GM]}}
                ['ґудзик']
            }
        }
    }
    & & \\ \\

    \theutterance \stepcounter{utterance}  
    & & & \multicolumn{2}{p{0.3\linewidth}}{
        \cellcolor[rgb]{0.9,0.9,0.9}{
            \makecell[{{p{\linewidth}}}]{
                \texttt{\tiny{[GM$|$GM]}}
                {'word': 'ґудзик', 'assignment': 'team'}
            }
        }
    }
    & & \\ \\

    \theutterance \stepcounter{utterance}  
    & & & \multicolumn{2}{p{0.3\linewidth}}{
        \cellcolor[rgb]{0.9,0.9,0.9}{
            \makecell[{{p{\linewidth}}}]{
                \texttt{\tiny{[GM$|$GM]}}
                {'word': 'ґудзик', 'assignment': 'team'}
            }
        }
    }
    & & \\ \\

    \theutterance \stepcounter{utterance}  
    & & & \multicolumn{2}{p{0.3\linewidth}}{
        \cellcolor[rgb]{0.9,0.9,0.9}{
            \makecell[{{p{\linewidth}}}]{
                \texttt{\tiny{[GM$|$GM]}}
                {'word': 'лимон', 'assignment': 'opponent'}
            }
        }
    }
    & & \\ \\

    \theutterance \stepcounter{utterance}  
    & & & \multicolumn{2}{p{0.3\linewidth}}{
        \cellcolor[rgb]{0.9,0.9,0.9}{
            \makecell[{{p{\linewidth}}}]{
                \texttt{\tiny{[GM$|$GM]}}
                {'hidden': {'team': ['маса', 'пляшка', 'собака'], 'innocent': ['ангел', 'коло', 'амазонка', 'концерт', 'вчитель', 'стовп', 'америка'], 'opponent': [], 'assassin': ['столиця']}, 'revealed': {'team': {'team': ['лук', 'день', 'англія', 'казино', 'акції', 'ґудзик'], 'innocent': [], 'opponent': ['муха'], 'assassin': []}, 'opponent': {'team': [], 'innocent': [], 'opponent': ['час', "зв'язок", 'двір', 'вентилятор', 'лінійка', "п'єса", 'лимон'], 'assassin': []}}}
            }
        }
    }
    & & \\ \\

    \theutterance \stepcounter{utterance}  
    & & \multicolumn{4}{p{0.6\linewidth}}{
        \cellcolor[rgb]{0.9,0.9,0.9}{
            \makecell[{{p{\linewidth}}}]{
                \texttt{\tiny{[P1$\langle$GM]}}
                \texttt{Слово ґудзик було відгадано правильно. Хід вашого товариша по команді закінчився.} \\
\texttt{Тепер надайте ще одну підказку, пов'язану з деякими з ваших командних слів, що залишилися, та список пов'язаних цільових слів. Пам'ятайте: почніть підказку з 'ПІДКАЗКА: ', перейдіть на новий рядок і почніть ваш список цільових слів, розділених комами, з 'ЦІЛІ: '. Зверніть увагу: деякі слова було вилучено зі списків порівняно з попередніми запитами.} \\
\\ 
\texttt{Ваші командні слова, що залишилися: маса, пляшка, собака.} \\
\texttt{Слова, що залишилися для вашого суперника: лимон.} \\
\texttt{Слова{-}відволікачі, що залишилися: ангел, коло, амазонка, концерт, вчитель, стовп, америка.} \\
\texttt{Слова{-}вбивці, що залишилися: столиця.} \\
            }
        }
    }
    & \\ \\

    \theutterance \stepcounter{utterance}  
    & \multicolumn{4}{p{0.6\linewidth}}{
        \cellcolor[rgb]{0.9,0.9,0.9}{
            \makecell[{{p{\linewidth}}}]{
                \texttt{\tiny{[P1$\rangle$GM]}}
                \texttt{ПІДКАЗКА: лабрадор} \\
\texttt{ЦІЛІ: собака, пляшка} \\
            }
        }
    }
    & & \\ \\

    \theutterance \stepcounter{utterance}  
    & & & \multicolumn{2}{p{0.3\linewidth}}{
        \cellcolor[rgb]{0.9,0.9,0.9}{
            \makecell[{{p{\linewidth}}}]{
                \texttt{\tiny{[GM$|$GM]}}
                \texttt{лабрадор} \\
            }
        }
    }
    & & \\ \\

    \theutterance \stepcounter{utterance}  
    & & & \multicolumn{2}{p{0.3\linewidth}}{
        \cellcolor[rgb]{0.9,0.9,0.9}{
            \makecell[{{p{\linewidth}}}]{
                \texttt{\tiny{[GM$|$GM]}}
                ['собака', 'пляшка']
            }
        }
    }
    & & \\ \\

    \theutterance \stepcounter{utterance}  
    & & & \multicolumn{2}{p{0.3\linewidth}}{
        \cellcolor[rgb]{0.9,0.9,0.9}{
            \makecell[{{p{\linewidth}}}]{
                \texttt{\tiny{[GM$|$GM]}}
                {'word': 'собака', 'assignment': 'team'}
            }
        }
    }
    & & \\ \\

    \theutterance \stepcounter{utterance}  
    & & & \multicolumn{2}{p{0.3\linewidth}}{
        \cellcolor[rgb]{0.9,0.9,0.9}{
            \makecell[{{p{\linewidth}}}]{
                \texttt{\tiny{[GM$|$GM]}}
                {'word': 'пляшка', 'assignment': 'team'}
            }
        }
    }
    & & \\ \\

    \theutterance \stepcounter{utterance}  
    & & & \multicolumn{2}{p{0.3\linewidth}}{
        \cellcolor[rgb]{0.9,0.9,0.9}{
            \makecell[{{p{\linewidth}}}]{
                \texttt{\tiny{[GM$|$GM]}}
                \texttt{opponent has won} \\
            }
        }
    }
    & & \\ \\

\end{supertabular}
}

\end{document}
